\section{Related Work}
\label{sec:related}

Our work sits at the intersection of three research lines: adversarial robustness benchmarking,
architecture comparison studies, and LLM trustworthiness evaluation. We survey each in turn,
identifying the limitations that motivate our contribution.

\subsection{Adversarial Robustness Benchmarks}

\citet{Ribeiro2020Beyond} introduced CheckList, a behavioral testing paradigm that evaluates
NLP models across a matrix of linguistic capabilities and perturbation types. CheckList shifts the
evaluation frame from aggregate accuracy to structured failure mode analysis, but does not compare
architecture families. Building on this, \citet{Wang2021Adversarial} released AdvGLUE --- a
multi-category adversarial benchmark derived from GLUE tasks using 14 attack methods covering
character-level noise, word-level substitution, and semantic perturbations. AdvGLUE is the primary
evaluation suite in our study. \citet{Nie2020Adversarial} developed ANLI (Adversarial NLI), where
examples are collected via human-model adversarial interaction, providing a benchmark resistant to
statistical artifacts from automatic perturbation methods. We use ANLI Round 3 (ANLI-R3) as the
human-crafted evaluation partition.

A key limitation of these benchmarks in prior use is their application for leaderboard-style scalar
comparison. We address this by treating AdvGLUE and ANLI-R3 as sources of per-category $\Delta^*$
observations and subjecting the resulting attack-category $\times$ model matrix to multivariate
analysis.

\subsection{Architecture Comparison Studies}

\citet{Neerudu2023Robustness} compare BERT, GPT-2, and T5 on GLUE perturbations and find that
GPT-2 (decoder-only) shows greater robustness than BERT (encoder-only) and T5 (encoder-decoder).
This directional finding is consistent with our $\eta^2=0.293$ result and the attention-topology
hypothesis. However, \citet{Neerudu2023Robustness} evaluate only three models (one per family),
report no effect size, apply no multivariate test, and do not use AdvGLUE or ANLI. Our work extends
their directional observation with formal effect-size quantification ($\eta^2=0.293$), a nine-model
pool, a six-category $\Delta^*$-vector representation, and permutation MANOVA with statistical
controls.

\citet{Zhang2026Robust} find that decoder-only models show lower explanation flip rates versus
encoder baselines in enterprise NLP, consistent with architecture-family robustness differences in
specific deployment contexts. \citet{Yang2024Assessing} compare LLaMA, OPT, and T5 under white-box
attacks and observe that model size and structure affect robustness, but again without formal
between-family effect-size quantification or $\Delta^*$-vector representation.

\subsection{LLM Trustworthiness Evaluation}

\citet{Sun2024TrustLLM} evaluate 16 LLMs across six trustworthiness dimensions using over 30
datasets. TrustLLM's robustness dimension covers adversarial perturbations relevant to our study,
but does not stratify by architecture family or compute per-attack-category $\Delta^*$-vectors.
Our work provides an architecture-family-stratified analysis of the robustness dimension that
TrustLLM evaluates at the model level. \citet{Otmakhova2025FLUKE} demonstrate that capability
does not imply robustness --- motivation consistent with our approach of measuring $\Delta^*$ as
distinct from clean accuracy. \citet{Cohen1988Statistical} provides the multivariate effect-size
conventions and power analysis framework we apply throughout.

\paragraph{Summary.} Prior work establishes directional architecture-family robustness differences
and documents them across diverse benchmarks. What is missing is a formal, effect-size-quantified,
scale-matched, multivariate analysis of between-family $\Delta^*$-vector clustering under
statistical controls --- and a validated reusable pipeline enabling community replication.
